\documentclass[12pt]{article}
\usepackage{amsmath,amssymb}

\begin{document}
\section*{{2019 USAMO Problems/Problem 5}}
Source: AoPS Wiki

\subsection*{Solution}
We claim that all odd $m, n$ work if $m+n$ is a positive power of 2. Proof:
We first prove that $m+n=2^k$ works. By weighted averages we have that $\frac{n(\frac{m}{n})+(2^k-n)\frac{n}{m}}{2^k}=\frac{m+n}{2^k}=1$ can be written, so the solution set does indeed work. We will now prove these are the only solutions. Assume that $m+n\ne 2^k$ , so then $m+n\equiv 0\pmod{p}$ for some odd prime $p$ . Then $m\equiv -n\pmod{p}$ , so $\frac{m}{n}\equiv \frac{n}{m}\equiv -1\pmod{p}$ . We see that the arithmetic mean is $\frac{-1+(-1)}{2}\equiv -1\pmod{p}$ and the harmonic mean is $\frac{2(-1)(-1)}{-1+(-1)}\equiv -1\pmod{p}$ , so if 1 can be written then $1\equiv -1\pmod{p}$ and $2\equiv 0\pmod{p}$ which is obviously impossible, and we are done. -Stormersyle

\end{document}
